%
% frame-visibility-lifecycle_buggy.tex
%
% FIDELITY CONTROL for docs/frame-visibility-lifecycle.tex.  This is the
% Display's frame bookkeeping AS THE CODE STANDS -- close_frame that pops the
% frame and forgets its scenes, and upsert_scene_in_frame's is_new branch that
% clears the minimized flag, requests focus, and takes the tab.  It exists to
% be violated: a model that cannot reproduce the defects it guards against is
% too abstract to trust.
%
% It differs from frame-visibility-lifecycle.tex in exactly three schemas:
% PushNewFrame, PushNewScene and Close.  Sys, Init, PushRepeat, DisposeScene,
% DisposeFrame, Minimize, Raise, SelectTab and ConsumeFocus are identical.
%
% Type-check:  fuzz -t docs/frame-visibility-lifecycle_buggy.tex
% Model-check (each goal must be FOUND -- see Section 3):
%   probcli docs/frame-visibility-lifecycle_buggy.tex -model_check -nodead \
%       -goal "autoFocused /= {}" -p DEFAULT_SETSIZE 2
%   probcli docs/frame-visibility-lifecycle_buggy.tex -model_check -nodead \
%       -goal "tabStolen /= {}" -p DEFAULT_SETSIZE 2
%   probcli docs/frame-visibility-lifecycle_buggy.tex -model_check -nodead \
%       -goal "userClosed /\ frames /= userClosed" -p DEFAULT_SETSIZE 2
%   probcli docs/frame-visibility-lifecycle_buggy.tex -model_check -nodead \
%       -goal "card({f|f:FRAME & f:frames & f:userClosed & vis(f)=vOpen})>0" \
%       -p DEFAULT_SETSIZE 2
%

\documentclass[a4paper,10pt,fleqn]{article}
\usepackage[margin=1in]{geometry}
\usepackage{fuzz}
\usepackage[colorlinks=true,linkcolor=blue,citecolor=blue,urlcolor=blue]{hyperref}

\begin{document}

\title{Frame Visibility as the Code Stands: a Fidelity Control}
\author{The variant that reaches the bad states}
\date{August 2026}
\maketitle

% ============================================================
\section{Purpose}
% ============================================================

\texttt{docs/frame-visibility-lifecycle.tex} claims that four invariants hold
of the Display's frame bookkeeping once bead \texttt{lux-mxvy.8} is
implemented.  The claim is worth nothing unless the same model, carrying
today's behaviour instead, \emph{reaches} the states those invariants forbid.
This document is that control.

Three schemas differ, and they are the three places today's code conflates
content with visibility.

\begin{description}
  \item[\texttt{Close}.]  Today \texttt{SceneReplica.close\_frame} calls
    \texttt{FrameBook.pop\_frame} and then, for every scene the frame held,
    \texttt{forget\_scene} and \texttt{widget\_state.discard}.  The frame and
    its scene ids leave the replica outright: a user's visibility decision is
    written into the content axis, and destroyed there.
  \item[\texttt{PushNewFrame} and \texttt{PushNewScene}.]  Today
    \texttt{upsert\_scene\_in\_frame} takes the \texttt{is\_new} branch ---
    \texttt{frame.active\_tab = msg.id}, \texttt{frame.minimized = False},
    \texttt{request\_focus(frame.frame\_id)} --- whenever the scene id is not
    already in the frame.  A content event writes all three
    visibility-owned values.  This is DES-060's shipped ``announce on
    arrival'', which R8 of DES-065 retires.
\end{description}

Everything else --- the state schema, the initial state, the repeat push, both
disposals, minimize, raise, tab selection, and the focus consume --- is
character-for-character the specification's.  The ghosts
(\texttt{userClosed}, \texttt{autoFocused}, \texttt{tabStolen}) are modelling
devices with no counterpart in the code, and they are what let a state be
judged against what the user actually did; \texttt{Close} here therefore still
records the user's intent even though the code it models retains nothing.

% ============================================================
\section{The Model}
% ============================================================

\begin{zed}
  [FRAME, SCENE]
\end{zed}

\begin{zed}
  VIS ::= vOpen | vMinimized | vClosed
\end{zed}

\begin{schema}{Sys}
  frames : \power FRAME \\
  content : \power SCENE \\
  sceneFrame : SCENE \pfun FRAME \\
  vis : FRAME \pfun VIS \\
  activeTab : FRAME \pfun SCENE \\
  focus : \power FRAME \\
  userClosed : \power FRAME \\
  autoFocused : \power FRAME \\
  tabStolen : \power FRAME
\where
  \dom vis = frames \\
  \dom activeTab = frames \\
  \dom sceneFrame = content \\
  \ran sceneFrame \subseteq frames \\
  \forall f : frames @ f \in \ran sceneFrame \\
  \forall f : frames @ activeTab~f \in content \\
  \forall f : frames @ sceneFrame~(activeTab~f) = f \\
  focus \subseteq frames \\
  \# focus \leq 1
\end{schema}

\begin{schema}{Init}
  Sys'
\where
  frames' = \emptyset \\
  content' = \emptyset \\
  sceneFrame' = \emptyset \\
  vis' = \emptyset \\
  activeTab' = \emptyset \\
  focus' = \emptyset \\
  userClosed' = \emptyset \\
  autoFocused' = \emptyset \\
  tabStolen' = \emptyset
\end{schema}

% -------------------------------------------------------------------
\subsection{Push a New Scene into a New Frame --- DIFFERS}
% -------------------------------------------------------------------

The \texttt{is\_new} branch fires on the frame it has just created: focus is
requested, and \texttt{autoFocused} records that a content event did it.

\begin{schema}{PushNewFrame}
  \Delta Sys \\
  s? : SCENE \\
  f? : FRAME
\where
  s? \notin content \\
  f? \notin frames \\
  frames' = frames \cup \{ f? \} \\
  content' = content \cup \{ s? \} \\
  sceneFrame' = sceneFrame \cup \{ s? \mapsto f? \} \\
  vis' = vis \cup \{ f? \mapsto vOpen \} \\
  activeTab' = activeTab \cup \{ f? \mapsto s? \} \\
  focus' = \{ f? \} \\
  autoFocused' = autoFocused \cup \{ f? \} \\
  userClosed' = userClosed \\
  tabStolen' = tabStolen
\end{schema}

% -------------------------------------------------------------------
\subsection{Push a New Scene into an Existing Frame --- DIFFERS}
% -------------------------------------------------------------------

The same branch on a frame that already exists, whatever the user had made of
it: the minimized flag is cleared, focus is requested, and the arriving scene
takes the tab.

\begin{schema}{PushNewScene}
  \Delta Sys \\
  s? : SCENE \\
  f? : FRAME
\where
  s? \notin content \\
  f? \in frames \\
  content' = content \cup \{ s? \} \\
  sceneFrame' = sceneFrame \cup \{ s? \mapsto f? \} \\
  vis' = vis \oplus \{ f? \mapsto vOpen \} \\
  activeTab' = activeTab \oplus \{ f? \mapsto s? \} \\
  focus' = \{ f? \} \\
  autoFocused' = autoFocused \cup \{ f? \} \\
  tabStolen' = tabStolen \cup \{ f? \} \\
  frames' = frames \\
  userClosed' = userClosed
\end{schema}

% -------------------------------------------------------------------
\subsection{Push a Scene Already Held --- identical}
% -------------------------------------------------------------------

\begin{schema}{PushRepeat}
  \Xi Sys \\
  s? : SCENE
\where
  s? \in content
\end{schema}

% -------------------------------------------------------------------
\subsection{Dispose a Scene --- identical}
% -------------------------------------------------------------------

\begin{schema}{DisposeScene}
  \Delta Sys \\
  s? : SCENE \\
  f? : FRAME \\
  keep? : SCENE
\where
  s? \in content \\
  sceneFrame~s? = f? \\
  keep? \in \dom (sceneFrame \rres \{ f? \}) \\
  keep? \neq s? \\
  content' = content \setminus \{ s? \} \\
  sceneFrame' = \{ s? \} \ndres sceneFrame \\
  activeTab' = activeTab \oplus \{ f? \mapsto keep? \} \\
  frames' = frames \\
  vis' = vis \\
  focus' = focus \\
  userClosed' = userClosed \\
  autoFocused' = autoFocused \\
  tabStolen' = tabStolen
\end{schema}

\begin{schema}{DisposeFrame}
  \Delta Sys \\
  s? : SCENE \\
  f? : FRAME
\where
  s? \in content \\
  sceneFrame~s? = f? \\
  \dom (sceneFrame \rres \{ f? \}) = \{ s? \} \\
  content' = content \setminus \{ s? \} \\
  sceneFrame' = \{ s? \} \ndres sceneFrame \\
  frames' = frames \setminus \{ f? \} \\
  vis' = \{ f? \} \ndres vis \\
  activeTab' = \{ f? \} \ndres activeTab \\
  focus' = focus \setminus \{ f? \} \\
  userClosed' = userClosed \setminus \{ f? \} \\
  autoFocused' = autoFocused \setminus \{ f? \} \\
  tabStolen' = tabStolen \setminus \{ f? \}
\end{schema}

% -------------------------------------------------------------------
\subsection{Minimize --- identical}
% -------------------------------------------------------------------

\begin{schema}{Minimize}
  \Delta Sys \\
  f? : FRAME
\where
  f? \in frames \\
  vis~f? = vOpen \\
  vis' = vis \oplus \{ f? \mapsto vMinimized \} \\
  focus' = focus \setminus \{ f? \} \\
  frames' = frames \\
  content' = content \\
  sceneFrame' = sceneFrame \\
  activeTab' = activeTab \\
  userClosed' = userClosed \\
  autoFocused' = autoFocused \\
  tabStolen' = tabStolen
\end{schema}

% -------------------------------------------------------------------
\subsection{Close --- DIFFERS}
% -------------------------------------------------------------------

Today's \texttt{close\_frame}: the frame is popped and every scene it held is
forgotten.  \texttt{vClosed} is never assigned, because today it is not a
state --- closed \emph{is} the absence of the frame.  The scene ids return to
the unseen pool, which is what makes the next push read as new.

The ghost \texttt{userClosed} still records that the user did this.  It has no
counterpart in the code being modelled; that is precisely the point, since the
defect is that the code retains nothing.

\begin{schema}{Close}
  \Delta Sys \\
  f? : FRAME
\where
  f? \in frames \\
  vis~f? \in \{ vOpen, vMinimized \} \\
  frames' = frames \setminus \{ f? \} \\
  content' = content \setminus \dom (sceneFrame \rres \{ f? \}) \\
  sceneFrame' = \dom (sceneFrame \rres \{ f? \}) \ndres sceneFrame \\
  vis' = \{ f? \} \ndres vis \\
  activeTab' = \{ f? \} \ndres activeTab \\
  focus' = focus \setminus \{ f? \} \\
  userClosed' = userClosed \cup \{ f? \} \\
  autoFocused' = autoFocused \setminus \{ f? \} \\
  tabStolen' = tabStolen \setminus \{ f? \}
\end{schema}

% -------------------------------------------------------------------
\subsection{Raise --- identical}
% -------------------------------------------------------------------

\begin{schema}{Raise}
  \Delta Sys \\
  f? : FRAME
\where
  f? \in frames \\
  vis' = vis \oplus \{ f? \mapsto vOpen \} \\
  focus' = \{ f? \} \\
  userClosed' = userClosed \setminus \{ f? \} \\
  frames' = frames \\
  content' = content \\
  sceneFrame' = sceneFrame \\
  activeTab' = activeTab \\
  autoFocused' = autoFocused \\
  tabStolen' = tabStolen
\end{schema}

% -------------------------------------------------------------------
\subsection{Select a Tab --- identical}
% -------------------------------------------------------------------

\begin{schema}{SelectTab}
  \Delta Sys \\
  f? : FRAME \\
  s? : SCENE
\where
  f? \in frames \\
  s? \in \dom (sceneFrame \rres \{ f? \}) \\
  activeTab' = activeTab \oplus \{ f? \mapsto s? \} \\
  frames' = frames \\
  content' = content \\
  sceneFrame' = sceneFrame \\
  vis' = vis \\
  focus' = focus \\
  userClosed' = userClosed \\
  autoFocused' = autoFocused \\
  tabStolen' = tabStolen
\end{schema}

% -------------------------------------------------------------------
\subsection{Consume the Focus Request --- identical}
% -------------------------------------------------------------------

\begin{schema}{ConsumeFocus}
  \Delta Sys \\
  f? : FRAME
\where
  f? \in focus \\
  focus' = \emptyset \\
  frames' = frames \\
  content' = content \\
  sceneFrame' = sceneFrame \\
  vis' = vis \\
  activeTab' = activeTab \\
  userClosed' = userClosed \\
  autoFocused' = autoFocused \\
  tabStolen' = tabStolen
\end{schema}

% ============================================================
\section{The States This Variant Reaches}
% ============================================================

Each goal below is the negation of one of the specification's invariants, and
each must be reported \emph{found} here.  The specification reports every one
of them \emph{not found}; that difference is the evidence the model sees the
defects the code is being changed to remove.

\paragraph{Goal 1 --- content took focus (Invariant 1 violated).}
\begin{verbatim}
  probcli docs/frame-visibility-lifecycle_buggy.tex -model_check -nodead \
    -goal "autoFocused /= {}" -p DEFAULT_SETSIZE 2
\end{verbatim}
One step: \texttt{PushNewFrame}$(s_1, f_1)$ leaves \texttt{autoFocused} $=
\{f_1\}$.  This is DES-025's auto-focus and DES-060's announce-on-arrival,
both of which R8 retires.

\paragraph{Goal 2 --- content stole the tab (Invariant 2 violated).}
\begin{verbatim}
  probcli docs/frame-visibility-lifecycle_buggy.tex -model_check -nodead \
    -goal "tabStolen /= {}" -p DEFAULT_SETSIZE 2
\end{verbatim}
Two steps: \texttt{PushNewFrame}$(s_1, f_1)$ then
\texttt{PushNewScene}$(s_2, f_1)$.  The user is reading $s_1$; $s_2$ arrives
and the selection moves under them.

\paragraph{Goal 3 --- bug A: the closed frame cannot be raised (Invariant 3
violated).}
\begin{verbatim}
  probcli docs/frame-visibility-lifecycle_buggy.tex -model_check -nodead \
    -goal "userClosed /\ frames /= userClosed" -p DEFAULT_SETSIZE 2
\end{verbatim}
Two steps: \texttt{PushNewFrame}$(s_1, f_1)$ then \texttt{Close}$(f_1)$.
Afterwards \texttt{userClosed} $= \{f_1\}$ and \texttt{frames} $= \emptyset$,
so \texttt{Raise}$(f_1)$ is not enabled --- the precondition $f_1 \in frames$
fails.  In the running system this is \texttt{raise\_frame} answering
\texttt{raised: false} because \texttt{FrameBook} popped the frame, and it is
why an applet's menu-entry click cannot bring back a frame the user closed.

Note what this means for the two halves of the fix: retiring the
\texttt{is\_new} side effect \emph{without} making \texttt{Close} keep the
frame would leave the user with no way back at all.  Today the accidental
recreation is the only reopen path there is.

\paragraph{Goal 4 --- bug B: the closed frame reopens on a content push
(Invariant 4 violated).}
\begin{verbatim}
  probcli docs/frame-visibility-lifecycle_buggy.tex -model_check -nodead \
    -goal "card({f|f:FRAME & f:frames & f:userClosed & vis(f)=vOpen})>0" \
    -p DEFAULT_SETSIZE 2
\end{verbatim}
Three steps: \texttt{PushNewFrame}$(s_1, f_1)$, \texttt{Close}$(f_1)$,
\texttt{PushNewFrame}$(s_1, f_1)$.  The third step is enabled only because the
second returned $s_1$ to the unseen pool; the frame comes back
\texttt{vOpen}, focused, while \texttt{userClosed} still holds $f_1$.  In the
running system this is voxd's next track change reopening the music window
the user shut.

\paragraph{Status of the ProB leg.}  As recorded in the specification's
Verification section, \texttt{probcli} is not installed on the development
host --- \texttt{\$(HOME)/Applications/ProB/probcli} holds only a 314-byte
\texttt{probcli.zip} that is not a valid archive.  This control is
\texttt{fuzz}-clean and its four traces are hand-derived from the operation
schemas above.  The reachability searches themselves have not been run.

\end{document}
